\subsection{Ablation Study and Additional Variants}
\label{sec:ablation}

We evaluate ablation and extension variants to test whether the final two-gate
allocator can be improved by additional routing, extra-action, clustering, parsing,
or pool-miss logic.
Table~\ref{tab:ablation} summarizes the evaluated variants; implementation labels are retained in the table for reproducibility.
All variants were evaluated via postrun replay on existing candidate artifacts;
no additional inference calls were issued for ablation purposes.
None of the additional variants improved on Failure-Trace Allocator (FTA) with sufficient held-out evidence, and they are
reported to delimit the supported claim.

\medskip\noindent\textbf{Retained gate components.}
Low-Depth Guard (LDG) alone achieves 80.00\% on Aggregate-720 ($+2.36$~pp over the frontier
baseline), confirming that the Low-Depth Guard carries most of the combined gain.
External-Consensus Fallback (ECO) alone achieves 75.00\% on final-300---below Low-Depth Guard (LDG)---but contributes an
additional $+0.69$~pp when combined with Low-Depth Guard (LDG), reflecting its role in recovering
cases where Low-Depth Guard (LDG) does not fire.
A score-absence calibration variant was evaluated but applied to
zero examples in the primary splits, contributing no accuracy change; it is
excluded from the combined policy.

\medskip\noindent\textbf{Additional variants not included in the final policy.}
A TALE-default router produced zero switches on seeds~41 and~61,
equalling TALE performance on both splits, and only a marginal net positive
switch on seed~71 ($+1$ example over TALE, final-300 78.67\% vs.\ 78.33\%).
This outcome is $-8.00$~pp below Failure-Trace Allocator (FTA) on the same split.
Note that the Aggregate-720 accuracy for this router (75.28\%) numerically coincides
with that of the frontier baseline; this reflects the zero-switch behavior on
seeds~41 and~61 (the router reverts to TALE, and TALE plus the single seed-71 switch
happen to sum to the same total as the frontier), not an equivalence by design.
The TALE-default design assumed a reliable agreement region that did not
materialise consistently across evaluation seeds; see Section~\ref{sec:discussion}
for interpretation.
An extra-action variant showed a nonnegative pilot signal in the
initial 40-case overlapping evaluation, but an independent 80-case follow-up
yielded net $-1$ for extra frontier actions and net $-4$ for extra TALE
retries relative to Failure-Trace Allocator (FTA).
Early-stage signal that fails to replicate on disjoint data is not sufficient
for inclusion in the final policy.
A cluster-level selector achieved at most $+0.67$~pp in the best rule
variant on final-300, below the threshold for inclusion in the main method.
A robust parser variant produced zero accuracy-changing switches on
both final-300 and Aggregate-720, indicating that parse-failure residuals are
rare in this evaluation setting.
A pool-miss diagnostic found 23 of 24 all-methods-wrong residuals to be
generation-limited, with no selector-only remedy available; see the failure
analysis in Section~\ref{sec:failure_analysis}.
